% Reusable caption for the parental-weighting teaching figure.
% Requires graphicx. Define \Description if your document class does not provide it.
\begin{figure*}[!t]
\centering
\includegraphics[width=0.95\textwidth]{parental_weighting.pdf}
\caption{The same parents receive different weights when we describe fitness and when we predict selection. Each point marks the breeding values of one simulated parent for two traits. We show the same 600 points in all three maps and use all 20,000 parents for the calculations (one two-trait scenario from the Gaussian validation experiment for Figure 2 of the main analysis). Coordinates are in standard deviations from the starting mean. The contours join trait combinations with equal log expected reproductive output; log output is set to zero at the starting mean on each surface. (A) Each parent counts once. Parents below $-1$ standard deviation for trait~1 (purple) make up 15.9\% of the population. (B) The landscape and coordinates are unchanged. Dot area is now proportional to expected offspring number on the curved surface. The purple parents account for 28.5\% of expected output. (C) A plane fitted to log expected output on the same cloud assigns them 34.1\%. (D) Weighting all parents gives the mean marked by the diamond for the curved surface or the triangle for the plane. Distances from the starting mean (open circle) are 0.64 and 0.83 in standardised trait coordinates. Although the plane explains 90.2\% of the variation in log expected output, it overestimates the distance the parental mean shifts by 30\% in this example.}
\Description{Three contour maps show the same cloud of breeding values. The left map gives each individual equal weight. In the middle and right maps, dots get larger or smaller according to expected offspring number, first on the curved landscape, then on its fitted plane. In the magnified view, both means are displaced towards lower values of trait one and higher values of trait two. The plane predicts the larger displacement.}
\label{fig:tutorial-parental-weighting}
\end{figure*}
